\chapter{Firmware de control: ESP32-S3 y CYD}
\label{cap:firmware-control}

Este capítulo recorre las dos placas que convierten un dedo apoyado sobre un
cristal en una forma de onda lista para sonar. El reparto de tareas entre
procesadores, la elección del medio de cada enlace y el tamaño de las tramas ya
quedaron justificados en el capítulo~\ref{cap:diseno}, en los
apartados~\ref{sec:arquitectura} y~\ref{sec:enlaces}. Aquí interesa el cómo, es
decir, cómo se lee y se calibra el táctil, cómo se localiza la celda de la
rejilla y se mezclan sus vecinas, cómo se construyen las tramas y cómo se reduce
la onda para dibujarla. Todo lo que sigue corresponde a la fase de ejecución
descrita en el apartado~\ref{sec:fases}, en la que no interviene ninguna red
neuronal.

\section{Lectura y calibración del táctil}
\label{sec:tactil}

La \gls{cyd} es una placa genérica sin manual de fabricante, cuya documentación
es comunitaria \cite{CYDrepo}, así que los pines se confirmaron uno a uno contra
el esquemático que la acompaña. Integra una pantalla \gls{tft} de 240$\times$320
puntos y un táctil resistivo de cuatro hilos gobernado por un XPT2046, que es un
convertidor analógico-digital de aproximaciones sucesivas
(\gls{sar}) de 12 bits \cite{XPT2046}. De ahí sale el dato que gobierna toda la
calibración, y es que cada eje se lee como un entero entre 0 y 4095 en cuentas
del convertidor, sin ninguna relación directa con los píxeles de la pantalla.

Táctil y pantalla cuelgan de buses distintos, y no por casualidad. El táctil usa
el bus VSPI con un objeto propio y los pines de fábrica de la placa, mientras que
la pantalla se controla con la librería \texttt{TFT\_eSPI} \cite{TFTeSPI} sobre
el bus HSPI. Por defecto esa librería también tomaría el VSPI, con lo que dos
objetos acabarían disputándose el mismo periférico. Forzando el HSPI son dos
buses independientes que conviven sin estorbarse, y de paso los pines 12 a 15
son los nativos del HSPI en el ESP32, encaminados por la matriz \mbox{IO MUX}
\cite{EspressifESP32}, que es la ruta directa y rápida. El táctil se lee con la
librería de Stoffregen \cite{XPT2046Touchscreen}.

Esa configuración vive en el fichero de proyecto de PlatformIO y no en el
\texttt{User\_Setup.h} de la librería, que es donde normalmente se pone. El
motivo es práctico. Ese fichero se instala dentro del directorio de dependencias
descargadas, fuera del control de versiones, de modo que la configuración se
perdería al reinstalar la librería o al clonar el proyecto en otro ordenador.
Declarándola como definiciones de compilación viaja con el repositorio. Hay
además un detalle de la placa que conviene no reaprender, y es que la
retroiluminación arranca apagada. La pantalla puede estar pintando
perfectamente y verse negra, indistinguible de un cuelgue, así que se enciende a
mano en la primera línea del arranque. Justo después se dibuja un patrón de
barras de color con un marco, que permite validar el controlador, la orientación
y los colores sin depender todavía del ESP32-S3.

\subsection{De las cuatro esquinas a las dos ventanas}

La primera calibración se hizo tocando las cuatro esquinas del cristal con un
puntero, que al tener un área de contacto pequeña da lecturas más extremas y
limpias que el dedo. De ahí salieron los límites 133 a 3947 en el eje horizontal
y 193 a 3897 en el vertical.

Esos números dejaron de valer cuando la placa se montó girada 180 grados en el
panel, decisión de montaje que se explica en el capítulo~\ref{cap:hardware} y que
busca alejar el conector \gls{usb} de la pantalla \gls{oled}. Al girar la imagen
hay que girar también el táctil, y aquí apareció un hallazgo que ahorró trabajo.
La librería sí aplica la rotación a las lecturas crudas, restándolas de 4095, lo
que evita invertir el mapeo a mano. Pero, por eso mismo, los cuatro límites de
calibración quedan espejados y hay que rehacerlos cruzando mínimo con máximo,
con lo que pasan a 148 a 3962 y 198 a 3902.

La ventana impresa del panel se come parte del cristal, y esa parte hubo que
medirla. La predicción hecha sobre plano, que daba un recorte de 3,75~mm por el
lado izquierdo y nada en vertical, no resistió el banco. Con la placa ya montada,
tocando las cuatro esquinas del hueco, los límites reales salieron 450 a 3730 en
los dos ejes. El labio asimétrico del panel existe, pero pesan más el error de
montaje y un ligero giro del táctil respecto al hueco. Se intentó ajustar esquina
por esquina y se descartó, porque las dos lecturas superiores no coincidían entre
sí en el eje vertical, con unas 230 cuentas de desviación entre la izquierda y la
derecha. Esa discrepancia es la firma de un giro, y un mapeo rectangular no puede
corregir un giro, así que perseguir cuatro números que no cuadran entre ellos
solo habría movido el problema de sitio.

Quedaba un problema más difícil. El labio impreso apoya sobre el cristal en todo
el perímetro y genera pulsaciones fantasma pegadas al borde. Una de las medidas
dio la lectura cruda (3714, 643), a solo 16 cuentas del límite derecho, es decir
dentro del rectángulo que se había dado por útil. Ningún ajuste de esos límites
podía rechazarla sin rechazar también contacto legítimo del dedo. Se probó
entonces a filtrar por presión, aprovechando que el XPT2046 la estima y que el
plástico apoya de forma constante pero floja mientras el dedo aprieta más, y
también se descartó. La fiabilidad de esa lectura en esta placa concreta es
dudosa y, sobre todo, introducía un umbral que mal ajustado podía dejar el táctil
mudo.

\begin{figure}[htbp]
    \centering
    \includegraphics[width=0.70\textwidth]{ventanas.pdf}
    \caption{Las dos ventanas del táctil sobre el plano de cuentas crudas del
             convertidor. La de captura decide qué toques valen y es la que se
             mapea al rango de la coordenada; la visible solo sirve para acotar
             el área de dibujo. La pulsación fantasma medida cae dentro de la
             visible y fuera de la de captura.}
    \label{fig:ventanas}
\end{figure}

La solución fue dejar de usar un solo rectángulo y separar dos, recogidos en la
Figura~\ref{fig:ventanas}. La ventana \emph{visible}, de 450 a 3730, describe lo
que asoma por el hueco y solo se emplea para calcular el área de dibujo. La
ventana de \emph{captura}, de 600 a 3580, decide qué toques se dan por buenos y
deja 150 cuentas de margen por cada lado, unos 2,5~mm. El fantasma medido queda
fuera de ella con 134 cuentas de holgura.

El detalle que hace que esa holgura no cueste nada es cuál de las dos ventanas se
mapea al rango de la coordenada. Se mapea la de captura, no la visible, de modo
que el espacio latente se sigue alcanzando entero de extremo a extremo y lo único
que se pierde son esos 2,5~mm de recorrido del dedo por lado. Mapear la visible
habría dejado las esquinas del plano inalcanzables, que es justamente el defecto
que la recalibración venía a quitar. Y el rechazo es total, no un recorte. La
versión anterior sujetaba la lectura al rango con una función de saturación, así
que un contacto del plástico se enviaba igual, clavado contra el borde del
espacio latente. Ahora un toque fuera de la ventana de captura no genera envío ni
repinta nada, y el instrumento se queda exactamente como estaba.

El táctil se atiende cada 50~ms, sin bloquear el bucle, y solo se envía
coordenada cuando el punto se ha desplazado más de 300 unidades respecto al
último enviado. Como una celda de la rejilla ocupa 4369 unidades del rango de 16
bits, ese umbral es menos de la décima parte de una celda, suficientemente fino
para que el recorrido se perciba continuo y suficientemente grueso para no
saturar el enlace con la vibración de un dedo quieto.

\section{De la coordenada a la celda}
\label{sec:coordenada}

La \gls{cyd} normaliza la lectura al rango de 16 bits antes de enviarla, y esa
decisión tiene una consecuencia comprobada. Ni el giro de 180 grados ni la
recalibración de banco obligaron a tocar una sola línea del firmware del
ESP32-S3, porque este sigue recibiendo el mismo rango de siempre. Resolver el
giro en el procesador de control habría roto la separación de responsabilidades
del apartado~\ref{sec:arquitectura} y habría dejado a la \gls{cyd} enviando
coordenadas que no se corresponden con lo que ella misma muestra.

Del lado del ESP32-S3, la trama de seis bytes descrita en el
apartado~\ref{sec:enlaces} la decodifica una máquina de estados que busca el byte
de cabecera, lee los cuatro de datos y valida la suma de comprobación. Ante
cualquier fallo vuelve al estado inicial a buscar cabecera, de forma que un byte
perdido descuadra una trama pero no el flujo. Conviene señalar que esa entrada
convive con otra por \gls{usb}, en la que la coordenada se escribe a mano como
texto. No es un resto de depuración olvidado. Fue lo que permitió construir y
validar la interpolación antes de que existiera el firmware de la \gls{cyd}, y
es también lo que hace posibles las medidas del capítulo~\ref{cap:pruebas}, que
necesitan una coordenada conocida y estable.

\section{Interpolación bilineal sobre la rejilla}
\label{sec:bilineal}

Localizar la celda es una multiplicación. La coordenada de 16 bits se escala por
$15/65535$, con lo que se obtiene la posición real dentro de la rejilla de
$16\times16$. La parte entera de esa posición da la celda y la parte fraccionaria
da el peso hacia la vecina siguiente. Los dos ejes se tratan de forma
independiente, así que un punto asimétrico no recibe ningún trato especial.

El borde se resuelve sin caso particular. Cuando la coordenada vale el máximo, la
posición escalada da exactamente 15, la vecina de la derecha se sujeta sobre sí
misma y el peso hacia ella queda a cero. La operación sigue siendo la misma
combinación de cuatro tablas, solo que dos de ellas coinciden y aportan peso
nulo, con lo que no se sale del array ni hace falta una rama distinta.

Hay que respetar la convención de índices que fija la cabecera generada en el
apartado~\ref{sec:rejilla}, con la fila recorriendo el eje vertical y la columna
el horizontal. Es el detalle más fácil de equivocar de todo el firmware y el que
peor avisa cuando se equivoca, porque permutar los ejes no rompe nada visible.
La interpolación sigue devolviendo formas de onda perfectamente válidas, solo que
el espacio de timbres queda transpuesto bajo el dedo.

\begin{figure}[htbp]
    \centering
    \includegraphics[width=\textwidth]{bilineal.pdf}
    \caption{Interpolación bilineal para la coordenada (40000, 12000). A la
             izquierda, la celda alcanzada, los índices de las cuatro vecinas en
             la rejilla y el peso que recibe cada una. A la derecha, esas cuatro
             formas de onda y la que resulta de mezclarlas.}
    \label{fig:bilineal}
\end{figure}

La Figura~\ref{fig:bilineal} muestra el resultado para la coordenada con la que
después se validó el enlace. Cae en la celda que va de la columna 9 a la 10 y de
la fila 2 a la 3, con pesos de 0,155 y 0,747, de modo que la vecina inferior
izquierda se lleva 0,631 de la mezcla y la superior derecha apenas 0,039. Se
aprecia que la onda resultante no es una de las cuatro ni un promedio simple,
sino que sigue de cerca a la que más pesa.

Merece una aclaración cómo se hace la aritmética, porque el formato de
almacenamiento y el de cálculo no coinciden. Las tablas se guardan como enteros
de 16 bits en formato Q15, pero la mezcla se calcula en coma flotante de 32 bits
y solo se redondea al pasar el resultado a Q15, una vez por muestra. El ESP32-S3
tiene unidad de coma flotante en hardware, así que las mil y pico operaciones que
esto cuesta son despreciables frente a los 1,6~ms que después ocupa la
transferencia, y el cálculo se ejecuta a frecuencia de control, unas decenas de
veces por segundo como mucho. La ventaja de hacerlo así es que cada muestra sufre
un único redondeo en lugar de acumular el de cada paso intermedio. Tampoco hace
falta saturar el resultado, porque los cuatro pesos suman uno y están todos entre
cero y uno, de modo que la mezcla es una combinación convexa de cuatro valores
que ya caben en Q15 y no puede salirse del rango.

La rejilla se compila como una constante de $16\times16\times1024$ enteros de 16
bits, es decir los 512~kB del apartado~\ref{sec:parametros}, y vive en memoria de
programa. En el ESP32-S3 la caché mapea esa memoria al espacio de direcciones, de
modo que se lee como cualquier otro array, sin las macros de acceso que exigen
otras arquitecturas. En memoria de datos solo queda la tabla de trabajo, dos
kilobytes. La \gls{psram} de 8~MB de la placa no llega a intervenir en el camino
principal y queda como margen para la ampliación planteada en el
apartado~\ref{sec:alternativas}.

Para comprobar que esta implementación es correcta no basta con ver que produce
algo. Se escribió una reimplementación independiente en numpy que recalcula la
misma interpolación sobre el fichero de la rejilla y se compara muestra a muestra
con el volcado que el ESP32-S3 imprime por \gls{usb}. Es una comprobación de
equivalencia frente a una referencia, no de plausibilidad. El resultado se recoge
en el capítulo~\ref{cap:pruebas}.

\section{Envío de la tabla al procesador de audio}
\label{sec:emisor-spi}

El ESP32-S3 dispone de varios controladores \gls{spi}, pero dos de ellos quedan
dedicados a la memoria de programa y la \gls{psram} del propio encapsulado y no
son utilizables \cite{EspressifS3}. Se emplea el primero de los de propósito
general, sobre los \gls{gpio} 10 a 13. La elección tenía que esquivar además los
pines que ocupa la variante con memorias octales de esta placa, más el que ya
usaba la entrada del \gls{uart}, y esos cuatro eran los que quedaban libres sin
comprometer nada.

El envío es de maestro, a 10~MHz, en modo 0 y con el bit más significativo por
delante. El selector de chip se gobierna a mano en lugar de dejarlo al
periférico, porque la escritura en bloque de la librería no lo conmuta por su
cuenta. La transferencia de los 2054 bytes bloquea durante 1,6~ms, cosa que no
supone problema alguno a la cadencia de control con la que se llama.

El contenido de la trama es el del apartado~\ref{sec:enlaces}. Las 1024 muestras
se copian en bloque, sin recorrer el array muestra a muestra, porque los dos
extremos del enlace almacenan los enteros de 16 bits en el mismo orden de bytes
y no hay nada que reordenar. Sobre ese bloque se calcula el \gls{crc} de 16 bits
que cierra la trama.

Ese \gls{crc} está escrito tres veces de forma idéntica, en el emisor, en el
receptor y en el guion de validación. Duplicar código suele ser mala señal, pero
aquí es deliberado, ya que tener la misma función en los tres sitios permite
comparar el valor calculado en cada uno y distinguir un fallo de transporte de
un fallo de cálculo. La validación cruzada del enlace tiene una particularidad de
método que conviene contar. El primer intento se hizo deslizando el dedo por la
\gls{cyd}, y resultó inútil, porque con el dedo en movimiento no hay forma de
saber con certeza qué trama recibida corresponde a qué coordenada enviada. Se
repitió con una coordenada fija escrita por \gls{usb}, que es lo que da una
referencia inequívoca. Las cifras están en el capítulo~\ref{cap:pruebas}.

\section{El canal de vuelta y el dibujo de la onda}
\label{sec:pantalla}

Lo que el ESP32-S3 devuelve a la pantalla es dato de representación, nunca de
síntesis. La rejilla y la interpolación siguen viviendo solo en el procesador de
control, y la forma de onda real únicamente sale hacia el procesador de audio por
\gls{spi}. El enlace con la \gls{cyd} ya era bidireccional, así que abrir el
sentido de vuelta no costó más que añadir el hilo y un emisor.

\begin{figure}[htbp]
    \centering
    \includegraphics[width=\textwidth]{diezmado.pdf}
    \caption{Reducción de 1024 a 256 puntos para dibujar, sobre el nodo de la
             rejilla en el que más se separan los dos métodos. Promediar cada
             grupo de cuatro muestras y quedarse con una de cada cuatro dan
             prácticamente lo mismo en el grueso del ciclo, y difieren en los
             flancos.}
    \label{fig:diezmado}
\end{figure}

La reducción a 256 puntos de 8 bits, cuyo porqué ya se explicó en el
apartado~\ref{sec:enlaces}, no se hace tomando una muestra de cada cuatro sino
promediando cada grupo de cuatro antes de quedarse con un punto, y el paso de
Q15 a Q7 se resuelve con un desplazamiento de ocho bits.
La Figura~\ref{fig:diezmado} compara los dos métodos sobre el nodo de la rejilla
en el que más se separan. En el grueso del ciclo coinciden, y la diferencia se
concentra en los tramos de pendiente fuerte, donde llega a 68 cuentas de Q7 sobre
una amplitud máxima de 122. Es decir, sin promediar la figura de la pantalla
seguiría reconociéndose, pero colocaría mal los flancos, que es precisamente
donde el ojo juzga la forma de una onda.

Dos tropiezos de esta parte tienen causa identificada y merece la pena dejarlos
escritos. El primero fue que no llegaba ninguna trama buena, no que llegaran
corruptas. El buffer de recepción del puerto serie en esta plataforma es de 256
bytes por omisión, justo por debajo de los 260 que ocupa la trama, y hay que
agrandarlo antes de abrir el puerto. El segundo fue que las tramas se perdían
solo mientras el dedo se movía, que es justo cuando llegan. La causa era una
espera bloqueante en la lectura del táctil, que desatendía la recepción durante
ese tiempo; se sustituyó por una temporización no bloqueante.

El repintado se hace en dos pasadas, borrando entera la traza anterior con el
color de fondo antes de dibujar la nueva. Repintar el área de golpe con un
rectángulo relleno produce parpadeo en cada trama, y borrar e ir pintando de
forma intercalada se come trozos de la línea nueva allí donde las dos se cruzan.
El borrado arrastra también los puntos de la línea central de referencia, que se
repintan sobre la marcha.

Queda el encaje con el panel. Como este tapa un borde del cristal, dibujar sobre
la pantalla completa dejaría parte de la interfaz debajo del plástico, así que el
área útil se acotó a los píxeles que salen de convertir la ventana visible de la
Figura~\ref{fig:ventanas}. Dentro de ella se reparten un título, una banda de
estado con la coordenada, el número de secuencia y el contador de tramas
descartadas, y el recuadro de la onda. La primera versión de ese reparto dejaba
la banda de estado terminando en la misma fila en la que empezaba el recuadro, y
el texto se comía la línea superior del marco. El patrón de arranque cierra el
asunto de forma barata, porque dibuja un marco blanco justo en el límite del área
útil y basta con ver si los cuatro lados asoman por la ventana para saber si las
constantes cuadran con el hueco.

\section{Estado del subsistema}

Con esto queda cerrada la cadena que va del dedo a la trama de audio. La
\gls{cyd} lee el táctil, lo filtra contra la ventana de captura, lo normaliza y
lo envía; el ESP32-S3 localiza la celda, mezcla las cuatro vecinas y entrega la
tabla por \gls{spi}, además de devolver a la pantalla una versión reducida de lo
que está sonando. Las dos comprobaciones cuantitativas que sostienen este
capítulo, la equivalencia de la interpolación frente a la referencia y la
integridad del enlace, se recogen juntas en el capítulo~\ref{cap:pruebas} para
poder leerlas al lado del resto de medidas del sistema montado. Lo que llega al
otro extremo del bus son 2054 bytes cada vez que el dedo se mueve, y qué se hace
con ellos es el asunto del capítulo siguiente.
